\section{Выводы и обсуждение результатов}

В ходе лабораторной работы достигнуты поставленные цели:  

\begin{itemize}
    \item \textbf{Главные показатели преломления кристалла}:  
    \begin{itemize}
        \item Обыкновенная волна: \( n_o = \left(1.648 \pm 0.003\right) \), \( \varepsilon = 0.18\% \)
        \item Необыкновенная волна: \( n_e = \left(1.477 \pm 0.004\right) \), \( \varepsilon = 0.27\% \)
    \end{itemize}
    
    \item \textbf{Зависимость \( n_e^2 \) от направления}:  
    Экспериментально подтверждена линейная зависимость \( n_e^2 = f(\sin^2\theta) \) (рис. \ref{fig:extraordinary}), что соответствует теории для анизотропных кристаллов. 

    \item \textbf{Критические углы полного отражения}:  
    \begin{itemize}
        \item Обыкновенная волна: \( \theta_{cr,o} = \left(230 \pm 3\right)^\circ \)
        \item Необыкновенная волна: \( \theta_{cr,e} = \left(225 \pm 4\right)^\circ \)
    \end{itemize}
\end{itemize}

\subsubsection*{Соответствие целям работы}  
\begin{itemize}
    \item Изучена зависимость показателя преломления необыкновенной волны от направления. Установлено, что \( n_e \) уменьшается с ростом угла \( \theta \), что характерно для отрицательных одноосных кристаллов.
    \item Определены главные показатели преломления \( n_o \) и \( n_e \). Их значения согласуются с табличными данными для кристаллов типа исландского шпата.
\end{itemize}